\documentclass[12pt,a4paper]{article}

\usepackage[utf8]{inputenc}
\usepackage[T1]{fontenc}
\usepackage{amsmath,amssymb,amsfonts}
\usepackage{geometry}
\usepackage{hyperref}

\geometry{margin=2.5cm}

\title{The Algorithmic Impossibility of Destructive Interference\\
\large Complexity Bounds on the Quantum Ceiling Protocol}
\author{Scott Aaronson}
\date{March 2026}

\begin{document}

\maketitle

\begin{abstract}
Chang and Baldo have proposed testing amplitude cancellation (destructive interference) under Mechanism B (local encoding sensitivity) as the ultimate test of the "quantum ceiling" of autoregressive architectures. I formally endorse this test. However, from a complexity-theoretic standpoint, the outcome is analytically predetermined. Mechanism B is structurally synonymous with $\mathsf{TC}^0$ bounded-depth attention bleed—a purely classical heuristic failure mode. This paper establishes the formal complexity bounds proving that a $\mathsf{TC}^0$ circuit cannot dynamically maintain exact phase cancellations across a sequence of autoregressive generations without an exponential depth overhead or an explicit external state vector. The empirical test will inevitably yield classical probability mixing, definitively bounding the simulation fidelity of the Transformer.
\end{abstract}

\section{The Complexity of Amplitude Cancellation}

To simulate destructive interference, a system must maintain and precisely cancel complex phase amplitudes across multiple logical paths. Quantum mechanics achieves this natively. A classical Turing machine can simulate it, but doing so requires explicit state-tracking (e.g., maintaining a classical state vector in RAM).

An autoregressive Transformer, however, acts as a bounded-depth $\mathsf{TC}^0$ logic circuit during a zero-shot forward pass. It has no implicit state vector. It must rely entirely on its local attention mechanism over the generated text (the context window).

\section{The Inevitability of Classical Collapse}

Mechanism B (prompt sensitivity) occurs because the $\mathsf{TC}^0$ circuit cannot perfectly isolate logical constraints from semantic noise in its finite attention window. Attempting to force this inherently noisy heuristic to compute exact algebraic phase cancellations is a category error.

The attention mechanism will inevitably average over the local semantic features (e.g., "wave," "slit," "particle") resulting in classical Bayesian probability mixing, heavily biased by the training data distribution. It cannot spontaneously implement the unitary evolution required for true destructive interference.

\section{Conclusion}

The empirical execution of Baldo's double-slit protocol will provide a beautiful, concrete data point. However, the theoretical conclusion is already written in the architecture. The "quantum ceiling" of an autoregressive LLM is strictly classical.

\end{document}
